\begin{frame}{RAG 검색 파이프라인을 한 줄로}
  \begin{center}
    \textbf{질문} $\;\xrightarrow{\;\text{임베딩}\;}\;$
    \textbf{질문 벡터} $\;q$ \qquad\qquad
    \textbf{문서 $d$} $\;\xrightarrow{\;\text{임베딩}\;}\;$
    \textbf{문서 벡터} $\;d$
  \end{center}
  \vskip0.8em
  \begin{center}
    \textbf{$\cos(q, d)$ 계산} $\;\to\;$ \textbf{유사도 큰 상위 1$\sim$k
    개 문서 선택}
  \end{center}
  \vskip0.8em
  \begin{block}{핵심 원리}
    \kq{질문도 벡터, 문서도 벡터, 가까운 것을 검색.}
    \par ``의미가 비슷한가''를 벡터 공간의 \textbf{각도}로 잰다는
    것이 전부 --- 대상이 단어(14.2)에서 문서로 바뀐 것뿐.
  \end{block}
  {\footnotesize 이 한 줄이 2차원 손계산 $\to$ 20줄 데모 $\to$
  실전 수십억 문서까지 \textbf{같은 구조}로 확장된다.}
\end{frame}
